\begin{textsample}{POS Dim 1 – 2017 – Score 23.00 – 2017/t000541.txt}  \label{ex:f1_pos_056}
We know more about other \textbf{planets} than our own, and today, I want to show you a new type of robot designed to help us better understand our own \textbf{planet}.
It belongs to a category known in \textbf{the} oceanographic community as an unmanned \textbf{surface} vehicle, or USV.
And it uses no fuel.
Instead, it relies on wind power for propulsion.
And yet, it can sail around \textbf{the} globe for months at a time.
So I want to share with you why we built it, and what it means for you.
A few years ago, I was on a sailboat making its way across \textbf{the} Pacific, from San Francisco to Hawaii.
I had just spent \textbf{the} past 10 years working nonstop, developing video games for hundreds of millions of users, and I wanted to take a step back and look at \textbf{the} big picture and get some much-needed thinking time.
I was \textbf{the} navigator on board, and one evening, after a long session analyzing weather data and plotting our course, I came up on deck and saw this beautiful sunset.
And a thought \textbf{occurred} to me : How much do we really know about our oceans? \textbf{The} Pacific was stretching all around me as far as \textbf{the} eye could see, and \textbf{the} waves were \textbf{rocking} our boat forcefully, a sort of constant reminder of its untold power.
How much do we really know about our oceans?
I decided to find out.
What I quickly learned is that we don ' t know very much. \textbf{The} first reason is just how vast oceans are, covering 70 percent of \textbf{the} \textbf{planet}, and yet we know they drive complex planetary systems like global weather, which affect all of us on a daily basis, sometimes dramatically.
And yet, those activities are mostly invisible to us.
Ocean data is scarce by any standard.
Back on land, I had grown used to accessing lots of sensors — billions of them, actually.
But at \textbf{sea}, in situ data is scarce and expensive.
Why?
Because it relies on a small number of ships and buoys.
How small a number was actually a great surprise.
Our National Oceanic and Atmospheric Administration, better known as NOAA, only has 16 ships, and there are less than 200 buoys offshore globally.
It is easy to understand why : \textbf{the} oceans are an unforgiving place, and to collect in situ data, you need a big ship, capable of carrying a vast amount of fuel and large crews, costing hundreds of millions of dollars each, or, big buoys tethered to \textbf{the} ocean floor with a four-mile-long cable and weighted down by a set of train wheels, which is both dangerous to deploy and expensive to maintain.
What about satellites, you might ask?
Well, satellites are \textbf{fantastic}, and they have taught us so much about \textbf{the} big picture over \textbf{the} past few decades.
However, \textbf{the} problem with satellites is they can only see through one micron of \textbf{the} \textbf{surface} of \textbf{the} ocean.
They have relatively poor spatial and temporal resolution, and their signal needs to be corrected for cloud cover and land effects and other factors.
So what is going on in \textbf{the} oceans?
And what are we trying to measure?
And how could a robot be of any use?
Let ' s zoom in on a small cube in \textbf{the} ocean.
One of \textbf{the} key things we want to understand is \textbf{the} \textbf{surface}, because \textbf{the} \textbf{surface}, if you think about it, is \textbf{the} nexus of all air-sea interaction.
It is \textbf{the} interface through which all energy and gases must flow.
Our sun radiates energy, which is absorbed by oceans as heat and then partially released into \textbf{the} atmosphere.
Gases in our atmosphere like CO2 get dissolved into our oceans.
Actually, about 30 percent of all global CO2 gets absorbed.
Plankton and microorganisms release oxygen into \textbf{the} atmosphere, so much so that every other breath you take comes from \textbf{the} ocean.
Some of that heat generates evaporation, which creates clouds and then eventually leads to precipitation.
And pressure gradients create \textbf{surface} wind, which moves \textbf{the} moisture through \textbf{the} atmosphere.
Some of \textbf{the} heat radiates down into \textbf{the} deep ocean and gets stored in different layers, \textbf{the} ocean acting as some kind of planetary-scale boiler to store all that energy, which later might be released in short-term events like hurricanes or long-term phenomena like El Ni{\EmojiFont �}{\EmojiFont �}o.
These layers can get mixed up by vertical upwelling currents or \textbf{horizontal} currents, which are key in transporting heat from \textbf{the} tropics to \textbf{the} poles.
And of course, there is marine life, occupying \textbf{the} largest ecosystem in volume on \textbf{the} \textbf{planet}, from microorganisms to \textbf{fish} to marine mammals, like \textbf{seals}, dolphins and \textbf{whales}.
But all of these are mostly invisible to us. \textbf{The} challenge in studying those ocean variables at scale is one of energy, \textbf{the} energy that it takes to deploy sensors into \textbf{the} deep ocean.
And of course, many solutions have been tried — from wave-actuated \textbf{devices} to \textbf{surface} drifters to sun-powered \textbf{electrical} drives — each with their own compromises.
Our team breakthrough came from an unlikely source — \textbf{the} pursuit of \textbf{the} world speed record in a wind-powered land yacht.
It took 10 years of research and development to come up with a novel \textbf{wing} concept that only uses three watts of power to control and yet can propel a vehicle all around \textbf{the} globe with seemingly unlimited autonomy.
By adapting this \textbf{wing} concept into a marine vehicle, we had \textbf{the} genesis of an ocean drone.
Now, these are larger than they appear.
They are about 15 feet high, 23 feet long, seven feet deep.
Think of them as \textbf{surface} satellites.
They ' re laden with an \textbf{array} of science-grade sensors that measure all key variables, both oceanographic and atmospheric, and a live satellite link transmits this high- resolution data back to shore in real time.
Our team has been hard at work over \textbf{the} past few years, conducting missions in some of \textbf{the} toughest ocean conditions on \textbf{the} \textbf{planet}, from \textbf{the} Arctic to \textbf{the} tropical Pacific.
We have sailed all \textbf{the} way to \textbf{the} polar ice shelf.
We have sailed into Atlantic hurricanes.
We have rounded Cape Horn, and we have slalomed between \textbf{the} \textbf{oil} rigs of \textbf{the} Gulf of Mexico.
This is one tough robot.
Let me share with you recent work that we did around \textbf{the} Pribilof Islands.
This is a small group of \textbf{islands} deep in \textbf{the} cold Bering Sea between \textbf{the} US and Russia.
Now, \textbf{the} Bering Sea is \textbf{the} home of \textbf{the} walleye pollock, which is a whitefish you might not recognize, but you might likely have tasted if you \textbf{enjoy} \textbf{fish} sticks or surimi.
Yes, surimi looks like crabmeat, but it ' s actually pollock.
And \textbf{the} pollock fishery is \textbf{the} largest fishery in \textbf{the} nation, both in terms of value and volume — about 3. 1 billion \textbf{pounds} of \textbf{fish} caught every year.
So over \textbf{the} past few years, a fleet of ocean drones has been hard at work in \textbf{the} Bering Sea with \textbf{the} goal to help assess \textbf{the} size of \textbf{the} pollock \textbf{fish} stock.
This helps improve \textbf{the} quota system that ' s used to manage \textbf{the} fishery and help prevent a collapse of \textbf{the} \textbf{fish} stock and protects this fragile ecosystem.
Now, \textbf{the} drones survey \textbf{the} \textbf{fishing} ground using acoustics, i. e., a sonar.
This sends a sound wave downwards, and then \textbf{the} reflection, \textbf{the} echo from \textbf{the} sound wave from \textbf{the} seabed or schools of \textbf{fish}, gives us an idea of what ' s happening below \textbf{the} \textbf{surface}.
Our ocean drones are actually pretty good at this repetitive task, so they have been gridding \textbf{the} Bering Sea day in, day out.
Now, \textbf{the} Pribilof Islands are also \textbf{the} home of a large colony of fur \textbf{seals}.
In \textbf{the} 1950s, there were about two million individuals in that colony.
Sadly, these days, \textbf{the} population has rapidly declined.
There ' s less than 50 percent of that number left, and \textbf{the} population continues to fall rapidly.
So to understand why, our science partner at \textbf{the} National Marine Mammal Laboratory has fitted a GPS \textbf{tag} on some of \textbf{the} mother \textbf{seals}, glued to their furs.
And this \textbf{tag} measures location and depth and also has a really cool little camera that ' s triggered by sudden acceleration.
Here is a movie taken by an artistically inclined \textbf{seal}, giving us unprecedented insight into an \textbf{underwater} hunt deep in \textbf{the} Arctic, and \textbf{the} shot of this pollock prey just seconds before it gets devoured.
Now, doing work in \textbf{the} Arctic is very tough, even for a robot.
They had to survive a snowstorm in August and interferences from bystanders — that little spotted \textbf{seal} \textbf{enjoying} a ride.
Now, \textbf{the} \textbf{seal} \textbf{tags} have recorded over 200, 000 dives over \textbf{the} season, and upon a closer look, we get to see \textbf{the} individual \textbf{seal} tracks and \textbf{the} repetitive dives.
We are on our way to decode what is really happening over that foraging ground, and it ' s quite beautiful.
Once you superimpose \textbf{the} acoustic data collected by \textbf{the} drones, a picture starts to emerge.
As \textbf{the} \textbf{seals} leave \textbf{the} \textbf{islands} and swim from left to right, they are observed to dive at a relatively shallow depth of about 20 meters, which \textbf{the} drone identifies is populated by small young pollock with low calorific content. \textbf{The} \textbf{seals} then swim much greater distance and start to dive deeper to a place where \textbf{the} drone identifies larger, more adult pollock, which are more nutritious as \textbf{fish}.
Unfortunately, \textbf{the} calories expended by \textbf{the} mother \textbf{seals} to swim this extra distance don ' t leave them with enough energy to lactate their pups back on \textbf{the} \textbf{island}, leading to \textbf{the} population decline.
Further, \textbf{the} drones identify that \textbf{the} water temperature around \textbf{the} \textbf{island} has significantly warmed.
It might be one of \textbf{the} driving forces that ' s pushing \textbf{the} pollock north, and to spread in \textbf{search} of colder regions.
So \textbf{the} data analysis is ongoing, but already we can see that some of \textbf{the} pieces of \textbf{the} puzzle from \textbf{the} fur \textbf{seal} mystery are coming into focus.
But if you look back at \textbf{the} big picture, we are mammals, too.
And actually, \textbf{the} oceans provide up to 20 kilos of \textbf{fish} per human per year.
As we deplete our \textbf{fish} stocks, what can we humans learn from \textbf{the} fur \textbf{seal} story?
And beyond \textbf{fish}, \textbf{the} oceans affect all of us daily as they drive global weather systems, which affect things like global agricultural output or can lead to devastating destruction of lives and property through hurricanes, extreme heat and floods.
Our oceans are pretty much unexplored and undersampled, and today, we still know more about other \textbf{planets} than our own.
But if you divide this vast ocean in six-by-six-degree squares, each about 400 miles long, you ' d get about 1, 000 such squares.
So little by little, working with our partners, we are deploying one ocean drone in each of those boxes, \textbf{the} hope being that achieving planetary coverage will give us better insights into those planetary systems that affect humanity.
We have been using robots to study distant worlds in our solar system for a while now.
Now it is time to quantify our own \textbf{planet}, because we cannot fix what we cannot measure, and we cannot prepare for what we don ' t know.
Thank you.

% matched lemmas: array, device, electrical, enjoy, fantastic, fish, fishing, horizontal, island, occur, oil, planet, pound, rock, sea, seal, search, surface, tag, the, underwater, whale, wing
\end{textsample}
